\section{Conclusion}
\label{sec:conclusion}

The central finding of this paper is that checkpoint loading for LLM serving exhibits identifiable performance regimes: the optimal I/O backend changes systematically with checkpoint format, model scale, and deployment environment. This is not a single bandwidth number. It is a workload-shaped systems problem. The layout determines request geometry, the backend determines submission policy, runtime capability determines which interfaces can be used, and the integration layer determines how much of a backend improvement reaches users.

Tensora measures this path across two layouts, four backend policies, three observation layers, six public checkpoints, and three environments. The results reject a universal backend. H100 SafeTensors shifts from sync-favored small loads to \texttt{io\_uring}-favored large loads. H100 ServerlessLLM-style partitioning eliminates sync as a winner and rewards grouped range execution and planning. Python and vLLM reorder the details but preserve the practical consequence: loader choice changes public API time and serving initialization. H200 and A100 show that regime structure transfers better than backend identity.

The recommendation is therefore simple and bounded: load by design. Select from measured workload features, gate the choice by runtime capability, and expose explicit overrides for boundary cases. The right abstraction is not an I/O API name; it is a format-aware adaptive policy that recognizes checkpoint loading as a workload-dependent decision with measurable regime boundaries.
